% !TeX root = ../抽象代数笔记.tex
% 由主文档通过 \input 载入；请编译主文档。
\setcounter{chapter}{8} % 保留原书章号。
\chapter{群列、单群与可解群}\label{ch:solvable-groups}

\begin{description}
  \item[学习目标] 用逐层取商刻画群的复杂性，为根式可解性准备群论语言。
  \item[先修知识] 第\ref{ch:quotient-groups}、\ref{ch:actions}章；例题使用西罗定理。
\end{description}

\section{群列与单群}
约定次正规列（subnormal series）为有限链 $G=G_0\supseteq G_1\supseteq\cdots\supseteq G_r=\{e\}$，其中每个 $G_{i+1}\trianglelefteq G_i$；若每项都正规于 $G$，则称为正规列。定义因子群 $G_i/G_{i+1}$ 与单群（simple group），要求单群非平凡且没有非平凡真正规子群。用素数阶循环群建立最初例子。

\section{换位子与导出列}
固定换位子（commutator）$[a,b]=aba^{-1}b^{-1}$ 的约定，定义导出子群 $G'=[G,G]$ 为所有换位子生成的子群。证明 $G'$ 正规且 $G/G'$ 阿贝尔，并刻画其最小性：对 $N\trianglelefteq G$，商群 $G/N$ 阿贝尔当且仅当 $G'\subseteq N$。递归定义 $G^{(0)}=G$、$G^{(i+1)}=(G^{(i)})'$。

\section{可解群及其封闭性质}
定义可解群（solvable group）为导出列在有限步后达到平凡群的群，证明它等价于存在阿贝尔因子的次正规列。证明可解性对子群、商群及扩张封闭；扩张性质具体指对 $N\trianglelefteq G$，若 $N$ 与 $G/N$ 均可解，则 $G$ 可解。区分阿贝尔群、可解群与一般群。

\section{交错群与不可解性的第一个例子}
计算 $S_3$、$S_4$ 的导出列或构造阿贝尔因子的群列。通过 $A_5$ 的共轭类大小证明其为非阿贝尔单群，再推出 $A_5$ 与 $S_5$ 不可解。一般的 $A_n$（$n\geq5$）单性证明列为选学，并说明它在高次方程问题中的作用。

\section{合成列与练习规划}
定义合成列（composition series），将有限群合成列的存在性和若尔当--赫尔德定理（Jordan--Hölder theorem）列为选学：合成因子的同构类型及重数与合成列选择无关。核心练习要求计算换位子、检验一条群列是否满足条件，并给出可解但非阿贝尔的例子。

